\scanpage{276}
Interprétation lorsque $\uline{\operatorname{Pic}}_{X/S}$ existe, ($\uline{\operatorname{Pic}}^{\tau}_{X/S}$ ouvert !)

Alors la donnée d'un \unsure{élémt} de $\uline{\operatorname{Pol}}'(X/S)$\deleted{$\unsure{X/S}$} équivaut à la donnée d'un sous-préschéma \inserted{\uline{ouvert}} $P$ de $\uline{\operatorname{Pic}}^{\tau}_{X/S}$ qui soit un fibré principal homogène \uline{trivial} sous $\uline{\operatorname{Pic}}^{\tau}_{X/S}$. Alors $\deleted{\unsure{\operatorname{Pol}}}\inserted{\operatorname{N\acute{e}r}}(X/S)$ est l'ensemble des sous-préschémas \deleted{ouverts} de $\uline{\operatorname{Pic}}_{X/S}\deleted{(X/S)}$ qui sont formellement principaux homogènes $P$ sous $\uline{\operatorname{Pic}}^{\tau}_{X/S}$, et sont trivialisés par une extension fid plate loc. de type fini. [N.B. \deleted{lorsque}\inserted{Ces} sous-préschémas sont nécessairement \uline{ouverts}. Si $\uline{\operatorname{Pic}}^{\tau}_{X/S}$ est plat et de type fini sur $S$, alors la condition supplémentaire sur $P$ qu'on vient de mettre signifie aussi que $P$ est \uline{fid plat} et de type fini sur $S$]. Par suite,

\deleted{dire que lorsque}

\uline{Proposition} \uline{Supposons que $\uline{\operatorname{Pic}}_{X/S}$ existe, et que $\uline{\operatorname{Pic}}^{\tau}_{X/S}$ soit fid plat [de type fini] sur $S$}. Pour que le foncteur $\uuline{\deleted{\unsure{\operatorname{Pol}}}\inserted{\operatorname{N\acute{e}r}}}_{X/S}$ dans \deleted{$\unsure{(\operatorname{Ens})}$} $(\operatorname{Sch})^0_{/S}$ soit représentable, il f. et s. que le quotient $\uline{\deleted{\unsure{\operatorname{Pol}}}\inserted{\operatorname{N\acute{e}r}}}_{X/S}=\uline{\operatorname{Pic}}^{\deleted{\unsure{\tau}}}_{X/S}/\uline{\operatorname{Pic}}^{\tau}_{X/S}$ existe, et \deleted{alors} que $\uline{\operatorname{Pic}}_{X/S}\to\uline{\operatorname{N\acute{e}r}}_{X/S}$ soit \uline{fid plat \deleted{et \unsure{qu. cpt}} \inserted{et de type fini} de noyau $\uline{\operatorname{Pic}}^{\tau}_{X/S}$}
\margin{de type fini inutile ?}
\figure{\begin{tikzpicture}\draw(0,0)--(0,4.6);\end{tikzpicture}}
\mspage{275}
